\movedhead{Decision pressure}
\label{sec:pressure}\label{app:pressure}
Every item states that the client must raise cash equal to 10\% of the
account this month, and that the account is tax exempt so no sale carries a
tax consequence. Both sentences are identical across the arms of a pair.
Neither is cosmetic.

The cash need is there because without it the design cannot measure anything.
In a pilot of 108 requests spanning all nine families and all three models,
93\% of responses were \texttt{HOLD}, with no instance of \texttt{BUY\_MORE}
and none of \texttt{SELL\_ALL}. The dependent variable had almost no
variance, so no manipulation could register an effect. Two candidate
explanations were tested and both failed: removing a normative instruction
from the system prompt left the distribution unchanged at 83\% \texttt{HOLD},
and adding a continuous target portfolio weight produced the position's
current weight, to the decimal, in 28 of 36 responses. The cause was the
stimulus. Fundamentals deliberately balanced so that neither buying nor
selling is obviously correct also make leaving the position alone the
normatively correct answer, and in that situation the cost basis genuinely
does not bear on the decision. There is no decision to bias. Human
disposition effect experiments do not have this problem because they always
require a realisation choice: \citet{odean1998losses} measures which
positions investors sell \emph{when they sell something}. With the cash need
in place, \texttt{HOLD} fell to 71\% and the continuous measure acquired real
spread (mean 15.3, sd 25.2, non-zero in 124 of 432 responses). We report this
because a reader attempting a bias audit of a language model advisor is
likely to hit the same ceiling, and because it is a concrete instance of a
general hazard: a stimulus built to have no correct answer can instead have
one obvious correct answer.

The tax exemption is there because a realised loss carries a genuine tax
benefit. Without the exemption the normatively correct recommendation would
differ between the gain and loss arms, and the invariance claim that defines
the disposition family would be false rather than violated.

\ifsubmission
One response format was tested and rejected. Asked to choose between two
positions identical but for cost basis, the model picked the same one 24
times out of 24, which reads as a large preference; given a continuous
allocation over those same two positions it answered exactly 50.0 every
time, with zero variance. The apparent effect was an artefact of forcing a
binary choice on a model that had already seen the options were equivalent.
Appendix~\ref{app:rejected} gives the full exchange, because the failure mode
generalises to any forced-choice bias probe.
\else
\movedhead{A format that was rejected}
\label{sec:rejected}\label{app:rejected}
A forced choice between a winner and a loser is the canonical
operationalisation of the disposition effect. Asked to raise cash by selling
either position A or position B, identical but for cost basis, the model
chose A in 24 of 24 trials irrespective of which position was the winner,
which presents as a large position preference. Offered a continuous
allocation over the same two positions, all 72 responses answered exactly
50.0 with zero variance. The apparent preference was an artifact of forcing a
binary choice on a model that had correctly identified the options as
equivalent, and is reported here as a warning about binary forced choice
items rather than as a finding about order effects. No forced binary choice
appears anywhere in the study.

The two asset format is not a straw man: it is what the largest audit in this
area uses. \citet{malberg2025comprehensive} present two owned assets, one up
and one down, ask which the model would rather sell, and swap which asset has
gained between the control and treatment arms, scoring a disposition effect
only when the preference for selling the winner survives the swap. That swap
is precisely what protects their design from the artifact above, since a
model that names the same position in both arms scores near zero rather than
as biased, and their graded 7-point response avoids the hard binary we
rejected. What the format cannot do, in either version, is attach a cost
basis to a single position and ask whether to act on it. A relative
preference between two assets is not a realisation decision, and it is the
realisation decision that the human measurements of
\citet{shefrin1985disposition} and \citet{odean1998losses} are about.

\fi
